\section{Síntesis y ampliación}

En esta entrega se analizaron a fondo cuatro detalles técnicos clave de LightRAG:

\begin{enumerate}
    \item \textbf{La necesidad de KG-based RAG}: resuelve las limitaciones fundamentales de Naive RAG en el razonamiento multisalto y las consultas globales
    \item \textbf{Mecanismo de indexación incremental}: mediante fusión de descripciones al estilo MapReduce y compresión semántica, permite agregar documentos de forma dinámica
    \item \textbf{Visualización del KG}: ofrece una herramienta intuitiva de exploración del grafo basada en pyvis
    \item \textbf{Recuperación multimodo}: los tres modos Naive/Hybrid/Mixed cubren distintas necesidades de consulta
\end{enumerate}

\subsection{Lecturas adicionales}

\begin{itemize}
    \item El archivo \texttt{prompts.py} del código fuente de LightRAG: contiene todos los prompts clave (Entity Extraction, Summarization, etc.)
    \item Despliegue local de LangFuse: para monitorear las entradas y salidas de todas las llamadas a la API de LightRAG
    \item Microsoft GraphRAG: una solución de KG-based RAG similar a LightRAG, útil para un estudio comparativo
\end{itemize}

%% --- Fin del contenido --- %%

\end{document}
